\scanpage{358}
% @ head of sheet
\mspage{82}
% @ head, continuation of bracket
de la démonstration de 5.5 b) $\big]\!\big]$

% @ upper third, lines 2–6
Soit $x$ un objet (point) de $X$, $x_i$ son image dans $X_i$. Considérons (via \deleted{\unsure{la}}\unsure{ma|la} construction \unsure{standart|standard} des catégories des chemins) \deleted{\unsure{qu'on} peut itérer \unsure{fonct.}} les $\Omega^1(X,x)$, $\Omega^1(X_i,x_i)$. On trouve
\[
\Omega^1(X,x)\simeq\underrightarrow{\lim}_{I}\Omega^1(X_i,x_i),
\]
% @ centre, lines 7–11
et on vérifie que c'est là encore un système inductif $L_1$-$\unsure{\pi|\kappa}_0$-adapté. En itérant, on trouve pour $n\in\mathbb N^*$
\[
\Omega^n(X,x)\simeq\underrightarrow{\lim}_{I}\Omega^n(X_i,x_i),
\]
limite inductive $L_1$-$\pi_0$-adaptée. Utilisant 5.6, on en conclut que pour
% @ lower middle, lines 13–16
tout $n$, l'e\deleted{\unsure{ns}} sous-ens (supposons $I$ un ens. ordonné, que \deleted{\unsure{di}} \unsure{diable}, ça suffit !)
\[
I_n=\{i\in I\mid \Omega^n(X_i,x_i)\text{ est }0\text{-connexe}\}
\]
est cofinal dans $I$, de plus il est stable par $\underrightarrow{\lim}$ filtrantes de cardinal $\leq\pi_0$.
% @ foot, intersection and reference
Donc
\[
I_\infty\overset{\mathrm{def}}{=}\bigcap I_n
\]
est aussi cofinal dans $I$ (cf $\overline{\underline{\mathrm{XVIII}}}$ 13.10.2, \unsure{p.|n.|v.} 302).
% archive sheet number: 357
